\documentclass{utmthesis}
\usepackage{graphicx}
\usepackage{url} 
\usepackage[pages=some]{background}
\usepackage{algorithm,algorithmic}
\usepackage{algpseudocode}
\usepackage{amsmath}

\begin{document}

% Required information
\title{Uncertainty in Recurrent Neural Networks}
\author{Alireza Samar}
\degree{Master of Philosophy}
\specialization{Machine Learning}
\intakeyear{2016}
\faculty{Advanced Informatics School}
\titledate{April 2017}
\award{4}
% Options for Award 
% 1. Bachelor Degree Project Report
% 2. Master's Project Report (By course work)
% 3. Master's Dissertation (By course work and research)
% 4. Master's Thesis (By research)
% 5. Doctor of Philosophy Thesis
% 6. Other PhD Thesis
% 7. Generic PhD Thesis
% 8. Thesis Proposal
\superone{Dr. Siti Sophiayati Yuhaniz}
%\supertwo{M.Y. Other Supervisor}
%\superthree{Third SV}
%\superfour{Fourth SV}
%\superfive{Fifth SV}

% Option for two-page printing
\newgeometry{top=2.5cm,left=4cm,right=2.5cm,bottom=2.5cm,twoside}

% Option to add watermark page
% Comment for final version	
\backgroundsetup{scale=1,angle=0,opacity=.1,hshift=0.25in,vshift=-0.5in,contents={\includegraphics[width=5cm]{figs/utm02.jpg}}}
% \watermarkpage

% Mandatory pages
%\coverpage
%\superpage
%\certification
%\frontmatter
\maketitle
%\declaration

%\begin{dedication}
%Dedication\
%\end{dedication}

%\begin{acknowledgement}
%Acknowledgement
%\end{acknowledgement}

\begin{abstract}
Deep learning has outperformed in various fields from computer vision, and language processing to physics, biology, and manufacturing. This means the deep or multi-layer architecture of neural networks are being extensively used in these fields; for instance convolutional neural networks (CNN) as image processing tools, and recurrent neural networks (RNN) as sequence processing model.

However, in traditional sciences fields such as physics and biology, model uncertainty is crucial, especially in time series models where delay can’t be tolerated. In this work, I aim to propose a novel theoretical framework and develop tools to measure uncertainty estimates, especially in deep recurrent neural networks.

This work also tackles a widely known difficulty of training recurrent neural networks, vanishing gradient by proposing a novel architecture of RNN that compute weighted average unit on past iteration.

\end{abstract}

%\begin{abstrak}
%Ini adalah abstrak Bahasa Melayu
%\end{abstrak}

\tableofcontents
\listoftables
\listoffigures

%List of abbreviation 
\listofabbre
\addabbre{ANN}{Artificial Neural Network}
\addabbre{RNN}{Recurrent Neural Network}
\addabbre{BBB}{Bayes by Backprop}
\addabbre{CNN}{Convolutional Neural Network}
\addabbre{KL}{Kullback-Leibler}

%List of symbols 
\listofsymbols
\addsymbol{$\gamma$}{Whatever}
\addsymbol{$\sigma$}{Whatever}
\addsymbol{$\varepsilon$}{Whatever}

%List of appendices
%\listofappendices

\onehalfspacing
\mainmatter

\chapter{Literature Review}
\label{chap:lit.review}

\section{State-of-the-Arts}

\section{Limitations}
\begin{enumerate}
\item Mentor~Graphics 2
\begin{enumerate}
\item item 3
\end{enumerate}
\item item 4
\end{enumerate}

\section{Research Gaps}
The processing at layer-5%
\footnote{In this thesis, OSI model is used.} is done ...

\chapter{Literature Review}
\label{chap:lit.review}

\section{State-of-the-Arts}

\section{Limitations}
\begin{enumerate}
\item Mentor~Graphics 2
\begin{enumerate}
\item item 3
\end{enumerate}
\item item 4
\end{enumerate}

\section{Research Gaps}
The processing at layer-5%
\footnote{In this thesis, OSI model is used.} is done ...

\chapter{Research Methodology}
\label{chap:method}

\section{Operational Framework}
This research work can be described in the operational framework as illustrated in \textbf{Figure \ref{fig:opfw}}. It contains a structured set of activities have been designed to accomplish the objectives. The operational framework also monitors the progress of the project in several steps.

\begin{figure}
	\centering
	\includegraphics[scale=0.25]{./figs/opfw}
	\caption[Research Operational Framework]{Research Operational Framework.}
	\label{fig:opfw}
\end{figure}

\section{Planning and Research}
In this phase, previous works and the state-of-the-art has been extensively studied. Similar works like \textit{line search}, \textit{dynamic evaluation}, and \textit{learning to optimise} have reviewed to figure out research gap as described vastly in \textbf{Chapter \ref{chap:lit.review}}.

\section{Investigation of Model Confidence}
The quality of uncertainty estimates and the performance of the model in general is being affected by the form of posterior in variational inference.

That being said, by adopting the posterior locally, or sharpening, to a batch, there is a chance to improve the performance using gradients upon the batch.

\section{Propose and Develop a Technique}
While previous works train a parametric model, the proposed approach treat these as free parameters (so that they can adapt more quickly to the non-stationary distribution with respect to parameters). In this work, gradient information has been used to inform a variational posterior to reduce the variance of Bayesian Neural Networks. Therefore, although it is similar in flavour, the underlying motivations are quite different.

This technique could gives flexible form  to the typical assumption of Gaussian posterior (GP), that variational inference is applied to neural networks to reduces variance.

\section{Evaluation: Language Modeling}
The experiments result of this work will be tested in Chapter 5, on the Penn Treebank \cite{Marcus1993} benchmark, a task consisting on next word prediction. Using the network architecture from \cite{Zaremba2014} as a baseline and for which there is an open source implementation\footnote{\url{https://github.com/tensorflow/models/blob/master/tutorials/rnn/ptb/ptb_word_lm.py}}.

\begin{table}[t]
	\caption{Word-level perplexity on the Penn Treebank language modelling task (lower is better).}
	\label{tab:ptb}
	\vskip 0.15in
	\centering
	\small
	\begin{tabular}{r|c|c}
		\textbf{Model}             & \textbf{Validation} & \textbf{Test} \\ 
		\hline
		\abovespace
		LSTM & 120.7 & 114.5\\
		LSTM dropout                        & 	86.2	        &	82.1 \\
		SGLD    	                        &	N/A				&	109.1		\\
		SGLD dropout	                    &	N/A				&	99.6		\\
		Variational LSTM (tied weights)     &	81.8        	&	79.7 \\
		Variational LSTM (tied weights, MS)	&	N/A		        &	79.0		\\
		\hline 
		& &  \\
		\shortstack{\textbf{Model} \\ \textbf{with Dynamic Evaluation}}
		& \textbf{Validation} & \textbf{Test} \\
		\hline
		\abovespace
		LSTM dropout                      & 	79.7	        &	77.1 \\
		\hline
	\end{tabular}
	\vskip -0.1in
\end{table}

A vast report on recent advances on the Penn Treebank benchmark \cite{Kim2016}, \cite{Zilly2016}, \cite{Merity2016} that shows achievement of perplexities of 70.9 on the test set.

\section{Chapter Summary}
This chapter the operational framework of this research work has been described in details of each step which navigate the project from planning to evaluation.
\chapter{Literature Review}
\label{chap:lit.review}

\section{State-of-the-Arts}

\section{Limitations}
\begin{enumerate}
\item Mentor~Graphics 2
\begin{enumerate}
\item item 3
\end{enumerate}
\item item 4
\end{enumerate}

\section{Research Gaps}
The processing at layer-5%
\footnote{In this thesis, OSI model is used.} is done ...

%\chapter{Literature Review}
\label{chap:lit.review}

\section{State-of-the-Arts}

\section{Limitations}
\begin{enumerate}
\item Mentor~Graphics 2
\begin{enumerate}
\item item 3
\end{enumerate}
\item item 4
\end{enumerate}

\section{Research Gaps}
The processing at layer-5%
\footnote{In this thesis, OSI model is used.} is done ...

%\chapter{Literature Review}
\label{chap:lit.review}

\section{State-of-the-Arts}

\section{Limitations}
\begin{enumerate}
\item Mentor~Graphics 2
\begin{enumerate}
\item item 3
\end{enumerate}
\item item 4
\end{enumerate}

\section{Research Gaps}
The processing at layer-5%
\footnote{In this thesis, OSI model is used.} is done ...


\bibliographystyle{utmthesis-numbering}
\bibliography{reference}

%\appendix
%\chapter{Do not use long titles.}
%\chapter{Pseudo-codes}
%\chapter{Time-series Results}

%This is required to make List of Appendices possible. Remove when have no appendix.
%\endmatter
\end{document}
